% !TEX root = main.tex

We present the evaluation in five parts: standalone multiplier
comparison, in-core instance analysis, full-core hybrid vs.\ baseline
comparison, updated full-core switching-case analysis from benchmark VCD
replay, and benchmark-level dependency-distance analysis.

%% ================================================================
\subsection{Standalone Multiplier Comparison}\label{sec:eval:standalone}

Table~\ref{tab:standalone_power} reports the post-place-and-route
power breakdown for each multiplier synthesised and routed in
isolation under identical testbench stimuli.

\begin{table}[t]
  \caption{Standalone multiplier power and area (NanGate 45\,nm,
           100\,MHz, post-P\&R).}
  \label{tab:standalone_power}
  \centering
  \small
  \begin{tabular}{lrrr}
    \toprule
    \textbf{Metric} & \textbf{\textsc{mul}} & \textbf{\textsc{mule}} & $\boldsymbol{\Delta}$ \\
    \midrule
    Internal power (mW) & 0.6515 & 0.1383 & $-78.8$\% \\
    Switching power (mW)& 0.5601 & 0.0878 & $-84.3$\% \\
    Leakage power (mW)  & 0.0895 & 0.0392 & $-56.2$\% \\
    \textbf{Total power (mW)} & \textbf{1.301} & \textbf{0.265} & $\boldsymbol{-79.6}$\textbf{\%} \\
    \midrule
    Dynamic power (mW)  & 1.2116 & 0.2261 & $-81.3$\% \\
    Leakage fraction     & 6.88\% & 14.77\% & --- \\
    \midrule
    Cell count           & 2\,302  & 822    & $-64.3$\% \\
    Cell area ($\mu$m$^{2}$) & 4\,514.6 & 2\,135.4 & $-52.7$\% \\
    \bottomrule
  \end{tabular}
\end{table}

\textsc{mule} consumes \textbf{79.6\% less total power} than
\textsc{mul}, with reductions across all three components.  The
switching power reduction (84.3\%) is the largest contributor,
reflecting the dramatically smaller combinational cone.
\textsc{mule}'s leakage fraction (14.77\%) is higher than \textsc{mul}'s
(6.88\%) in relative terms because its dynamic power is so low;
in absolute terms, its leakage is still 56.2\% lower owing to the
52.7\% area reduction.

%% ================================================================
\subsection{In-Core Instance Analysis}\label{sec:eval:incore}

Table~\ref{tab:incore_power} presents the same power decomposition for
the multiplier instances measured \emph{within} the hybrid core during
the \texttt{tb\_mul\_compare} workload. Both units are present in the
same simulation and same netlist; the differences arise purely from
their micro-architectural activity.

\begin{table}[t]
  \caption{In-core multiplier instance power (Hybrid configuration,
           post-P\&R, VCD-annotated).}
  \label{tab:incore_power}
  \centering
  \small
  \begin{tabular}{lrrr}
    \toprule
    \textbf{Metric} & \textbf{u\_mul} & \textbf{u\_mule} & $\boldsymbol{\Delta}$ \\
    \midrule
    Internal power (mW)  & 0.5954 & 0.0871 & $-85.4$\% \\
    Switching power (mW) & 0.4663 & 0.0280 & $-94.0$\% \\
    Leakage power (mW)   & 0.1263 & 0.0399 & $-68.4$\% \\
    \textbf{Total power (mW)} & \textbf{1.188} & \textbf{0.155} & $\boldsymbol{-87.0}$\textbf{\%} \\
    \midrule
    \% of core power (8.270\,mW) & 14.36\% & 1.87\% & --- \\
    Dynamic power (mW)   & 1.0617 & 0.1151 & $-89.2$\% \\
    \bottomrule
  \end{tabular}
\end{table}

The in-core reduction is even larger than standalone (87.0\%
vs.\ 79.6\%) because the surrounding core logic effectively power-gates
\textsc{mule}'s inputs when no multiply is in flight, reducing
switching power by 94.0\%.  The \textsc{mul} instance accounts for
14.36\% of total core power; \textsc{mule} accounts for only 1.87\%.

%% ================================================================
\subsection{Energy Per Instruction and Energy--Delay Product}%
\label{sec:eval:epi}

Table~\ref{tab:epi_edp} consolidates EPI and EDP for both multipliers
across standalone and in-core contexts, with and without leakage.

\begin{table}[t]
  \caption{EPI and EDP comparison (standalone and in-core).}
  \label{tab:epi_edp}
  \centering
  \small
  \begin{tabular}{lrrrr}
    \toprule
    & \multicolumn{2}{c}{\textbf{EPI (pJ)}} & \multicolumn{2}{c}{\textbf{EDP (aJ$\cdot$s)}} \\
    \cmidrule(lr){2-3}\cmidrule(lr){4-5}
    \textbf{Context} & \textsc{mul} & \textsc{mule} & \textsc{mul} & \textsc{mule} \\
    \midrule
    Standalone (total)   & 26.02 & 13.25 & 520.4 & 662.5 \\
    Standalone (dynamic) & 24.23 & 11.31 & 484.6 & 565.5 \\
    In-core (total)      & 23.76 &  7.75 & 475.2 & 387.3 \\
    In-core (dynamic)    & 21.23 &  5.76 & 424.6 & 288.0 \\
    \bottomrule
  \end{tabular}
\end{table}

Two observations stand out.
\emph{First}, \textsc{mule} wins in EPI across all contexts: even in
the standalone case, where back-to-back exercise maximises its average
power, the 5-cycle latency is overcome by the $\sim\!5\times$ power
reduction, yielding a 49.1\% lower EPI.  In-core, the reduction widens
to \textbf{67.4\%} because idle-period power gating further suppresses
\textsc{mule}'s switching activity.

\emph{Second}, the EDP picture depends on context.  Standalone,
\textsc{mul} wins EDP because \textsc{mule}'s longer latency is
penalised quadratically.  \emph{In-core}, however, \textsc{mule}'s
dramatically lower power overcomes the latency penalty, yielding a
\textbf{32.2\% EDP improvement} on dynamic power---the most physically
meaningful comparison, as it isolates activity-dependent energy.

%% ================================================================
\subsection{Hybrid vs.\ Baseline Core Comparison}\label{sec:eval:hybrid}

Table~\ref{tab:core_power} compares total core power, multiplier
power, and area for the Hybrid (MUL\,+\,MULE) and Baseline
($2\!\times\!$MUL) configurations.

\begin{table}[t]
  \caption{Core-level comparison: Hybrid vs.\ $2\!\times\!$MUL
           baseline (post-P\&R, VCD-annotated).}
  \label{tab:core_power}
  \centering
  \small
  \begin{tabular}{lrrr}
    \toprule
    \textbf{Metric} & $\boldsymbol{2\!\times\!}$\textbf{MUL} & \textbf{Hybrid} & $\boldsymbol{\Delta}$ \\
    \midrule
    \multicolumn{4}{l}{\emph{Core power}} \\
    \quad Internal (mW)  & 4.532 & 4.430 & $-2.3$\% \\
    \quad Switching (mW) & 2.532 & 2.650 & $+4.7$\% \\
    \quad Leakage (mW)   & 1.230 & 1.190 & $-3.3$\% \\
    \quad \textbf{Total (mW)} & \textbf{8.295} & \textbf{8.270} & $\boldsymbol{-0.30}$\textbf{\%} \\
    \midrule
    \multicolumn{4}{l}{\emph{Multiplier power}} \\
    \quad u\_mul (mW)     & 2.020 & 1.188 & $-41.2$\% \\
    \quad u\_mul2 / u\_mule (mW) & 1.017 & 0.155 & $-84.8$\% \\
    \quad \textbf{Total (mW)} & \textbf{3.037} & \textbf{1.343} & $\boldsymbol{-55.8}$\textbf{\%} \\
    \midrule
    \multicolumn{4}{l}{\emph{Area}} \\
    \quad Core area ($\mu$m$^{2}$)  & 62\,640.9 & 61\,009.5 & $-2.6$\% \\
    \quad Core cell count            & 29\,153   & 28\,302   & $-2.9$\% \\
    \bottomrule
  \end{tabular}
\end{table}

Although total core power is nearly identical ($-0.30$\%), the internal
composition changes substantially. The combined multiplier power drops
by \textbf{55.8\%}, from 3.037\,mW (36.6\% of core) to 1.343\,mW
(16.2\% of core). This freed power budget can be allocated to other
micro-architectural features or to voltage/frequency reduction.
The core area decreases by 2.6\% because \textsc{mule} occupies only
2\,133.9\,$\mu$m$^{2}$ versus 4\,515.9\,$\mu$m$^{2}$ for the second
\textsc{mul}.

The slight increase in switching power ($+4.7$\%) in the hybrid core
is attributable to different routing topologies produced by the placer
when the second multiplier's footprint changes; it does not offset the
internal and leakage reductions.

%% ================================================================
\subsection{Updated Full-Core Switching Cases (2026)}
\label{sec:eval:switching2026}

To align the mechanism-level discussion with the latest validated
full-core benchmark resources, Table~\ref{tab:switching_cases_update}
summarizes the explicit latency-hidden switching workloads measured
through matched VCD replay. The key distinction is between
\emph{architecture-level} comparison (switched Hybrid vs.\ switched
2$\times$MUL) and \emph{rewrite-level} comparison (switched Hybrid
vs.\ baseline Hybrid).

To avoid baseline mixing, we report savings in three lanes throughout this
section. Lane A (architecture-level) answers whether Hybrid is better than the
homogeneous 2$\times$MUL core under the same workload and implementation flow;
this lane captures the larger power-oriented advantage discussed by the
break-even model. Lane B (rewrite-level) answers whether adding
\texttt{mul}$\rightarrow$\texttt{mule} substitutions improves or worsens the
Hybrid baseline on that same benchmark. Lane C (strict realized) applies the
benchmark-level exact-free acceptance rule and reports non-zero savings only
when total cycles are preserved exactly; timing-only headroom is shown
separately as potential, not realized gain.

\begin{table*}[t]
\centering
\caption{Updated switching outcomes from full-core benchmark VCD replay
(NanGate 45\,nm, 100\,MHz). Classification is based on switched Hybrid
E/op relative to switched 2$\times$MUL E/op.}
\label{tab:switching_cases_update}
\small
\begin{tabular}{l r r r r r l}
\toprule
Benchmark & Switched MUL sites & Hybrid power $\Delta$ vs 2$\times$MUL (\%) & Hybrid cycles $\Delta$ vs 2$\times$MUL (\%) & Hybrid E/op $\Delta$ vs 2$\times$MUL (\%) & Switched Hybrid E/op $\Delta$ vs Baseline Hybrid (\%) & Class \\
\midrule
\texttt{fft\_butterfly\_latency\_hidden\_mule} & 3 & -10.452162 & +0.294277 & -10.188643 & +2.775525 & Advantageous \\
\texttt{complex\_mul\_vec\_latency\_hidden\_mule} & 3 & -10.575752 & +1.816408 & -8.951443 & +5.199038 & Advantageous \\
\texttt{outer\_product\_latency\_hidden\_mule} & 1 & -9.990279 & +5.860363 & -4.715382 & +14.147057 & Advantageous \\
\texttt{matmul\_tiled\_latency\_hidden\_mule} & 1 & -9.222178 & +10.928167 & +0.698174 & +17.020176 & Marginally disadvantageous \\
\texttt{fir\_unrolled\_latency\_hidden\_mule} & 4 & -10.207159 & +15.468864 & +3.682773 & +47.578747 & Disadvantageous \\
\texttt{fir\_direct\_latency\_hidden\_mule} & 1 & -11.511374 & +26.351994 & +11.807144 & +17.726848 & Strongly disadvantageous \\
\bottomrule
\end{tabular}
\end{table*}

Table~\ref{tab:switching_cases_update} is therefore the right place to state
the \emph{architecture-level} energy result: on the best favorable switched
workloads, Hybrid reduces full-core E/op by about 9--10\% relative to the
homogeneous 2$\times$MUL baseline, and even the weakest favorable case still
improves E/op by 4.7\%. That claim should not be conflated with a
benchmark-level rewrite claim, however, because the same rows remain worse than
their own unswitched Hybrid baseline in the rewrite-level column. In other
words, the Hybrid architecture can be energy-favorable even when the current
binary rewriting policy does not yet produce an additional end-to-end win
inside Hybrid itself.

Three observations are important for the paper narrative.
\emph{First}, Hybrid consistently reduces average core power even in
negative switching cases. \emph{Second}, energy per operation is governed
by both power and runtime; when cycle inflation is large
(e.g., \texttt{fir\_direct\_latency\_hidden\_mule}), E/op degrades despite
lower power. \emph{Third}, in this dataset, every switched-Hybrid case is
worse than its own unswitched Hybrid baseline, so switching itself should
be presented as a conditional optimization rather than a universal gain.

This updated evidence is consistent with the framework of
Section~\ref{sec:method:framework}: local slack can make an iterative path
beneficial, but benchmark-level composition effects determine whether that
benefit survives at full-core scope.

%% ================================================================
\subsection{12-Benchmark In-House Exact-Free Shortlist}\label{sec:eval:1ghz_exactfree}

The earlier three-row replay table was not methodologically stable for
the current workspace-only rebuild: its switched rows already showed
benchmark-level cycle inflation, so they should not have been presented
as a strict exact-free subset. We therefore rebuild
Sections~\ref{sec:eval:1ghz_exactfree}
and~\ref{sec:eval:hybrid_mul_only_exactfree} around the reproducible
in-house benchmark family currently available in
\texttt{riscv\_dependency\_analysis}.

Before any switched replay, we first ask which kernels already hide the
architectural Hybrid overhead at whole-benchmark scope. Existing
baseline artifacts show that \texttt{horner} and
\texttt{unrolled\_dot} are the strongest current exact-free seeds. Both
keep identical total cycle counts across \texttt{core\_2standard},
\texttt{core\_hybrid}, and \texttt{core\_singlemul}, while still
retiring thousands of integer multiplies. That makes them better strict
exact-free candidates than the earlier replay rows
(\texttt{complex\_mul\_vec}, \texttt{fft\_butterfly},
\texttt{fir\_unrolled}), whose switched Hybrid runs already violate the
benchmark-level acceptance rule.

\begin{table*}[t]
\centering
\caption{Current in-house exact-free shortlist from the reproducible 12-source suite. Baseline cycle neutrality across \texttt{core\_2standard}, \texttt{core\_hybrid}, and \texttt{core\_singlemul} is used as the first screen before any switched 1\,GHz replay is attempted.}
\label{tab:exactfree_1ghz}
\small
\begin{tabular}{lrrrrl}
\toprule
Benchmark & Cycles ($2\!\times\!$MUL) & Cycles (Hybrid) & Cycles (singlemul) & Integer multiplies & Screening outcome \\
\midrule
\texttt{horner} & 114763 & 114763 & 114763 & 2048 & Shortlisted; first strict exact-free rerun target \\
\texttt{unrolled\_dot} & 29778 & 29778 & 29778 & 4096 & Shortlisted; first strict exact-free rerun target \\
\bottomrule
\end{tabular}
\end{table*}

The implication for replay is narrow but important. A switched benchmark
should only be replayed at 1\,GHz after it survives the same strict
benchmark-level acceptance test at 100\,MHz on this same in-house
suite. Under that rule, the current reproducible path does not justify
any new 1\,GHz switched E/op rows yet; the correct next replay queue is
\texttt{horner} and \texttt{unrolled\_dot}, not the earlier ad hoc
three-kernel subset.

\subsection{Strict Benchmark-Level Exact-Free Status on the In-House Suite}
\label{sec:eval:hybrid_mul_only_exactfree}

Table~\ref{tab:hybrid_mul_only_vs_latency_hidden_100mhz} reports the
current strict status of the multiply-screened portion of the same
in-house suite. Six kernels already had completed switched artifacts
and all six were rejected because switched Hybrid cycles exceeded the
Hybrid MUL-only baseline. We then reran the two strongest new shortlist
candidates. \texttt{horner} is now a measured rejection as well:
thresholds 1 and 2 both patch one static \texttt{mul} site, but cycles
still rise from 114763 to 125004, and threshold 3 finds no eligible
latency-hidden site. \texttt{unrolled\_dot} is also fully rejected by
the same strict rule: thresholds 1 and 2 patch four sites and raise
cycles from 29778 to 54344, while threshold 3 patches three sites and
still raises cycles to 46157. That leaves no measured strict exact-free
win yet in the current workspace-only suite.

\begin{table*}[t]
\centering
\caption{Current strict benchmark-level exact-free status on the multiply-screened in-house suite. The same 12-source benchmark family is used to define the shortlist and to report strict acceptance status.}
\label{tab:hybrid_mul_only_vs_latency_hidden_100mhz}
\small
\begin{tabular}{llll}
\toprule
Benchmark & Current evidence & Strict status & Notes \\
\midrule
\texttt{horner} & baseline 114763; switched 125004 at thresholds 1 and 2 & Rejected & Threshold 3 exposes no eligible latency-hidden site \\
\texttt{unrolled\_dot} & baseline 29778; switched 54344 (thr. 1/2), 46157 (thr. 3) & Rejected & Even the narrower threshold-3 rewrite still inflates cycles \\
\texttt{complex\_mul\_vec} & switched cycles 114798 vs baseline 109170 & Rejected & Cycle increase under current switched artifact \\
\texttt{fft\_butterfly} & switched cycles 87249 vs baseline 84951 & Rejected & Cycle increase under current switched artifact \\
\texttt{fir\_direct} & switched cycles 386651 vs baseline 322139 & Rejected & Cycle increase under current switched artifact \\
\texttt{fir\_unrolled} & switched cycles 214443 vs baseline 141372 & Rejected & Cycle increase under current switched artifact \\
\texttt{matmul\_tiled} & switched cycles 665234 vs baseline 562835 & Rejected & Cycle increase under current switched artifact \\
\texttt{outer\_product} & switched cycles 83256 vs baseline 71737 & Rejected & Cycle increase under current switched artifact \\
\texttt{matmul} & baseline artifact available & Unresolved & Switched artifact still missing \\
\texttt{dotprod} & partial artifacts only & Unresolved & Baseline/switched comparison still incomplete \\
\bottomrule
\end{tabular}
\end{table*}

This reframing makes the claim boundary explicit. The current workspace
already supports the architecture-level result that Hybrid can deliver
about 10\% whole-benchmark energy savings relative to switched
$2\!\times\!$MUL on selected kernels, but it does not yet support a
measured strict exact-free savings claim for the rewrite itself. Even
the two best current shortlist kernels fail the strict acceptance rule
under the first selective reruns completed so far. Any future 1\,GHz
replay row should therefore come only from a later accepted in-house row,
not from the presently rejected shortlist.

Taken together,
Tables~\ref{tab:switching_cases_update}, \ref{tab:exactfree_1ghz}, and
\ref{tab:hybrid_mul_only_vs_latency_hidden_100mhz} support three scoped
claims. First, at the architecture level, three of the six measured
switched cases reduce whole-benchmark E/op by 4.7--10.2\% relative to
switched $2\!\times\!$MUL, with roughly 9--11\% lower power; this is the
sense in which the Hybrid design delivers an ``around 10\%'' benchmark
energy saving. Second, at the rewrite level, every completed
switched-Hybrid case remains worse than its own unswitched Hybrid
baseline, so the current \texttt{mul}$\rightarrow$\texttt{mule} rewrite
policy should be presented as conditional rather than universally
beneficial. Third, within the reproducible 12-source suite,
\texttt{horner} and \texttt{unrolled\_dot} were the first
benchmark-level exact-free seeds worth retesting, but the completed
reruns reported here still reject them under the strict acceptance rule.
The break-even analysis in Section~\ref{sec:eval:breakeven} therefore
complements, rather than contradicts, the strict tables: it explains
when the shortlisted rewrite should be expected to help, while the strict
table reports that no currently rerun in-house kernel has yet cleared
benchmark-level acceptance.

%% ================================================================
\subsection{Compiler-Side Exact-Free Sweep (Mechanism Study)}\label{sec:eval:bench}

This subsection serves a different purpose than
Section~\ref{sec:eval:switching2026}. The 2026 switching table provides
the main full-core benchmark evidence used for paper claims. The analysis
below is retained as a mechanism-level study of exact-free site discovery
and composability under selective binary patching. In other words, it
explains \emph{why} some sites are safe to route to \texttt{mule}, while
Section~\ref{sec:eval:switching2026} reports the end-to-end core outcome
on the updated benchmark set.

We evaluate the alternative-execution-path idea on the 24 integer
kernels in the bare-metal benchmark harness used throughout this
work. Each kernel is executed in two primary configurations: an
all-\texttt{mul} baseline and an all-\texttt{mule} rewrite where every
eligible plain multiply is replaced by the custom instruction. The
benchmark return code alone is not treated as a correctness check.
Instead, we first generate a golden manifest of the printed
scalar/checksum outputs and then require every subsequent run to match
that manifest before the testbench reports PASS. Aggregate
all-\texttt{mul}/all-\texttt{mule} comparisons are formed only over
matched completed primary runs.

For exact-free replacement analysis, a static \texttt{mul} site is first
prefiltered using the implemented latency gap ($\Delta L = 2$), which in
this core corresponds to requiring the first dynamic consumer to be at
least three retired instructions later. That heuristic is not the final
metric, however: a site is counted only after patching the concrete
static PC to \texttt{mule} and rerunning the full benchmark with the
requirement that total cycle count remain unchanged.

Methodologically, each benchmark is compiled once into the bare-metal
harness and then evaluated in three controlled stages under the same RTL
and testbench configuration. First, an all-\texttt{mul} baseline run
collects total cycles, retired-instruction counts, and the retire trace
for every dynamic multiply. Second, an all-\texttt{mule} run rewrites
every eligible static multiply site to expose the cost of forcing the
slow path everywhere. Third, selective patch runs start from the
baseline binary, patch concrete static PCs to \texttt{mule}, and accept
a site only when the rerun still matches the golden output manifest and
preserves total cycle count exactly. This procedure keeps correctness,
instruction mix, and timing under the same measured testbench rather
than inferring the result from static distance alone.

\begin{table}[t]
\centering
\caption{Verified benchmark corpus summary. Primary rows count only matched all-\texttt{mul}/all-\texttt{mule} runs whose reported outputs matched the generated golden manifest. Exact-free counts come from binary-patched re-simulations that preserve total cycle count exactly.}
\label{tab:benchmark-verified-overview}
\begin{tabular}{lr}
\hline
Metric & Value \\
\hline
Benchmarks analyzed & 24 \\
Matched primary run pairs & 24 \\
Unsupported benchmarks & none \\
Aggregate all-\texttt{mul} cycles & 474,199 \\
Aggregate all-\texttt{mule} cycles & 515,463 \\
Retired instructions & 185,931 \\
Retired integer multiplies & 16,818 \\
Retired integer-multiply share & 9.05\% \\
Benchmarks with exact-free rewrites & 14 \\
Exact-free dynamic multiplies & 2,385 \\
Exact-free share of plain \texttt{mul} & 16.68\% \\
Exact-free share of retired instructions & 1.28\% \\
Patched exact-free static PCs & 131 \\
Exact-free total cycle delta & 0 \\
\hline
\end{tabular}
\end{table}

Table~\ref{tab:benchmark-verified-overview} shows that the verified
primary sweep covers all 24 kernels with no unsupported cases. The
aggregate retired-instruction and retired-multiply counts are identical
between the all-\texttt{mul} and all-\texttt{mule} runs, confirming that
the opcode substitution changes latency and issue behaviour rather than
the dynamic instruction mix. The cycle total nevertheless rises from
474,199 to 515,463 cycles when every multiply is forced onto
\textsc{mule}, reflecting the serialised single-outstanding nature of
the alternative path.

\begin{table*}[t]
\centering
\caption{Verified primary all-\texttt{mul} versus all-\texttt{mule} comparison. The table reports dynamic retired integer-multiply count and share for each benchmark, while overall slowdown is still computed from the underlying whole-benchmark all-\texttt{mul} and all-\texttt{mule} cycle totals. The exact-free columns cross-reference the independently validated selective-rewrite run to show which kernels admit timing-preserving conversions.}
\label{tab:benchmark-primary-comparison}
\resizebox{\textwidth}{!}{%
\begin{tabular}{l r r r r r r r}
\hline
Benchmark & $N$ & $K$ & Dyn. int. mul & Int.-mul share & Overall slowdown & Exact-free / mul & Patched PCs \\
\hline
\texttt{extended\_estrin} & 128 & -- & 2,304 & 33.48\% & 57.99\% & 11.11\% & 2 \\
\texttt{unrolled\_dot} & 128 & -- & 128 & 17.09\% & 39.92\% & 0.00\% & 0 \\
\texttt{rdr\_like\_corr} & 128 & 16 & 1,792 & 14.13\% & 27.78\% & 25.00\% & 1 \\
\texttt{matmul} & 16 & -- & 4,096 & 10.22\% & 15.34\% & 0.00\% & 0 \\
\texttt{fir} & 128 & 16 & 1,792 & 11.28\% & 14.03\% & 0.00\% & 0 \\
\texttt{estrin\_2level} & 128 & -- & 640 & 12.35\% & 13.79\% & 40.00\% & 2 \\
\texttt{estrin\_poly} & 128 & -- & 640 & 10.44\% & 9.55\% & 40.00\% & 2 \\
\texttt{outer\_product} & 16 & -- & 256 & 4.33\% & 7.42\% & 0.00\% & 0 \\
\texttt{horner} & 128 & -- & 128 & 7.03\% & 7.06\% & 0.00\% & 0 \\
\texttt{dotprod} & 128 & -- & 256 & 8.73\% & 5.67\% & 0.00\% & 0 \\
\texttt{sliding\_correlation} & 128 & -- & 1,240 & 17.33\% & 4.17\% & 40.00\% & 2 \\
\texttt{delayed\_reduction} & 128 & -- & 256 & 5.40\% & 3.25\% & 0.00\% & 0 \\
\texttt{complex\_bank} & 128 & -- & 1,024 & 13.13\% & 2.98\% & 37.50\% & 6 \\
\texttt{reduction\_tree\_mul} & 128 & -- & 256 & 9.07\% & 2.43\% & 25.00\% & 2 \\
\texttt{grouped\_correlation\_bank} & 128 & -- & 480 & 10.85\% & 2.07\% & 37.50\% & 3 \\
\texttt{unrolled\_dot16} & 128 & -- & 256 & 6.50\% & 0.39\% & 68.75\% & 11 \\
\texttt{unrolled\_dot8} & 128 & -- & 256 & 7.59\% & 0.26\% & 62.50\% & 5 \\
\texttt{multi\_accumulator} & 128 & -- & 256 & 7.43\% & 0.25\% & 62.50\% & 5 \\
\texttt{conv2d} & 16 & -- & 0 & 0.00\% & 0.00\% & -- & 0 \\
\texttt{filter\_bank\_4out} & 128 & -- & 250 & 0.87\% & 0.00\% & 100.00\% & 1 \\
\texttt{fir\_multi} & 128 & 4 & 0 & 0.00\% & 0.00\% & -- & 0 \\
\texttt{hello} & -- & -- & 0 & 0.00\% & 0.00\% & -- & 0 \\
\texttt{unrolled\_dot32} & 128 & -- & 256 & 6.03\% & 0.00\% & 87.50\% & 28 \\
\texttt{unrolled\_dot64} & 128 & -- & 256 & 5.57\% & -0.03\% & 95.31\% & 61 \\
\hline
\end{tabular}%
}
\end{table*}

Table~\ref{tab:benchmark-primary-comparison} exposes two distinct
regimes. Workloads that route most multiplies through dependence-dense
or repeatedly serialized regions suffer badly when every site is
rewritten to \texttt{mule}: \texttt{extended\_estrin},
\texttt{unrolled\_dot}, and \texttt{rdr\_like\_corr} slow down by
58.0\%, 39.9\%, and 27.8\%, respectively. At the other end of the
spectrum, the most aggressively unrolled kernels are effectively
cycle-neutral under the all-\texttt{mule} experiment:
\texttt{unrolled\_dot32} is exactly neutral and
\texttt{unrolled\_dot64} is marginally faster, while
\texttt{multi\_accumulator}, \texttt{unrolled\_dot16}, and
\texttt{reduction\_tree\_mul} retain sizeable exact-free subsets with
only small positive overheads.

\begin{table*}[t]
\centering
\caption{Benchmarks with non-zero verified exact-free rewrites. A static \texttt{mul} PC is counted only after the heuristic candidate is patched to \texttt{mule} and the full benchmark re-simulation preserves total cycle count.}
\label{tab:benchmark-exact-free}
\resizebox{\textwidth}{!}{%
\begin{tabular}{l r r r r r r}
\hline
Benchmark & $N$ & $K$ & Exact-free dyn. mul & Exact-free / mul & Exact-free / insn & Patched PCs \\
\hline
\texttt{filter\_bank\_4out} & 128 & -- & 125 & 100.00\% & 0.44\% & 1 \\
\texttt{unrolled\_dot64} & 128 & -- & 122 & 95.31\% & 2.66\% & 61 \\
\texttt{unrolled\_dot32} & 128 & -- & 112 & 87.50\% & 2.64\% & 28 \\
\texttt{unrolled\_dot16} & 128 & -- & 88 & 68.75\% & 2.23\% & 11 \\
\texttt{multi\_accumulator} & 128 & -- & 80 & 62.50\% & 2.32\% & 5 \\
\texttt{unrolled\_dot8} & 128 & -- & 80 & 62.50\% & 2.37\% & 5 \\
\texttt{estrin\_2level} & 128 & -- & 256 & 40.00\% & 4.94\% & 2 \\
\texttt{estrin\_poly} & 128 & -- & 256 & 40.00\% & 4.17\% & 2 \\
\texttt{sliding\_correlation} & 128 & -- & 248 & 40.00\% & 3.47\% & 2 \\
\texttt{complex\_bank} & 128 & -- & 192 & 37.50\% & 2.46\% & 6 \\
\texttt{grouped\_correlation\_bank} & 128 & -- & 90 & 37.50\% & 2.04\% & 3 \\
\texttt{rdr\_like\_corr} & 128 & 16 & 448 & 25.00\% & 3.53\% & 1 \\
\texttt{reduction\_tree\_mul} & 128 & -- & 32 & 25.00\% & 1.13\% & 2 \\
\texttt{extended\_estrin} & 128 & -- & 256 & 11.11\% & 3.72\% & 2 \\
\hline
\end{tabular}%
}
\end{table*}

\begin{table*}[t]
\centering
\caption{Complete per-benchmark \texttt{mul}-hiding breakdown ordered by exact-free share. Dynamic plain \texttt{mul} executions are partitioned into three classes: \emph{free} (patched to \texttt{mule} and confirmed cycle-neutral by full re-simulation), \emph{partial} (heuristic trace candidate that did not survive the benchmark-level exact-cycle acceptance test), and \emph{no-hide} (never a heuristic candidate; recurrence-dominated). Three benchmarks with zero static integer multiplies are included for completeness. Overhead is all-\texttt{mule} vs.\ all-\texttt{mul} cycle delta.}
\label{tab:benchmark-hiding-breakdown}
\resizebox{\textwidth}{!}{%
\begin{tabular}{l r r r r r r r r r}
\hline
Benchmark & Mul\% & Dyn.\ mul & Free & Free\% & Partial & Partial\% & No-hide & No-hide\% & MULE ovhd \\
\hline
\texttt{filter\_bank\_4out} & 0.87\% & 125 & 125 & 100.0\% & 0 & 0.0\% & 0 & 0.0\% & +0.00\% \\
\texttt{unrolled\_dot64} & 5.57\% & 128 & 122 & 95.3\% & 6 & 4.7\% & 0 & 0.0\% & -0.03\% \\
\texttt{unrolled\_dot32} & 6.03\% & 128 & 112 & 87.5\% & 16 & 12.5\% & 0 & 0.0\% & +0.00\% \\
\texttt{unrolled\_dot16} & 6.50\% & 128 & 88 & 68.8\% & 40 & 31.2\% & 0 & 0.0\% & +0.39\% \\
\texttt{multi\_accumulator} & 7.43\% & 128 & 80 & 62.5\% & 48 & 37.5\% & 0 & 0.0\% & +0.25\% \\
\texttt{unrolled\_dot8} & 7.59\% & 128 & 80 & 62.5\% & 48 & 37.5\% & 0 & 0.0\% & +0.26\% \\
\texttt{estrin\_2level} & 12.35\% & 640 & 256 & 40.0\% & 128 & 20.0\% & 256 & 40.0\% & +13.79\% \\
\texttt{estrin\_poly} & 10.44\% & 640 & 256 & 40.0\% & 128 & 20.0\% & 256 & 40.0\% & +9.55\% \\
\texttt{sliding\_correlation} & 17.33\% & 620 & 248 & 40.0\% & 248 & 40.0\% & 124 & 20.0\% & +4.17\% \\
\texttt{complex\_bank} & 13.13\% & 512 & 192 & 37.5\% & 192 & 37.5\% & 128 & 25.0\% & +2.98\% \\
\texttt{grouped\_correlation\_bank} & 10.85\% & 240 & 90 & 37.5\% & 60 & 25.0\% & 90 & 37.5\% & +2.07\% \\
\texttt{rdr\_like\_corr} & 14.13\% & 1,792 & 448 & 25.0\% & 896 & 50.0\% & 448 & 25.0\% & +27.78\% \\
\texttt{reduction\_tree\_mul} & 9.07\% & 128 & 32 & 25.0\% & 80 & 62.5\% & 16 & 12.5\% & +2.43\% \\
\texttt{extended\_estrin} & 33.48\% & 2,304 & 256 & 11.1\% & 1,536 & 66.7\% & 512 & 22.2\% & +57.99\% \\
\texttt{delayed\_reduction} & 5.40\% & 128 & 0 & 0.0\% & 128 & 100.0\% & 0 & 0.0\% & +3.25\% \\
\texttt{unrolled\_dot} & 17.09\% & 128 & 0 & 0.0\% & 96 & 75.0\% & 32 & 25.0\% & +39.92\% \\
\texttt{dotprod} & 8.73\% & 128 & 0 & 0.0\% & 0 & 0.0\% & 128 & 100.0\% & +5.67\% \\
\texttt{fir} & 11.28\% & 1,792 & 0 & 0.0\% & 0 & 0.0\% & 1,792 & 100.0\% & +14.03\% \\
\texttt{horner} & 7.03\% & 128 & 0 & 0.0\% & 0 & 0.0\% & 128 & 100.0\% & +7.06\% \\
\texttt{matmul} & 10.22\% & 4,096 & 0 & 0.0\% & 0 & 0.0\% & 4,096 & 100.0\% & +15.34\% \\
\texttt{outer\_product} & 4.33\% & 256 & 0 & 0.0\% & 0 & 0.0\% & 256 & 100.0\% & +7.42\% \\
\texttt{conv2d} & 0.00\% & -- & -- & -- & -- & -- & -- & -- & +0.00\% \\
\texttt{fir\_multi} & 0.00\% & -- & -- & -- & -- & -- & -- & -- & +0.00\% \\
\texttt{hello} & 0.00\% & -- & -- & -- & -- & -- & -- & -- & +0.00\% \\
\hline
\textbf{Total} & \textbf{9.05\%} & \textbf{14,297} & \textbf{2,385} & \textbf{16.7\%} & \textbf{3,650} & \textbf{25.5\%} & \textbf{8,262} & \textbf{57.8\%} & \textbf{+8.70\%} \\
\hline
\end{tabular}%
}
\end{table*}

Under this exact-free criterion, 2,385 of 14,297 dynamic plain
\texttt{mul} executions (16.7\%) can be shifted to \texttt{mule} with no
cycle cost, spanning 131 static PCs across 14 kernels. In aggregate,
those exact-free conversions account for 1.28\% of all retired
instructions. The highest verified fractions occur in
\texttt{filter\_bank\_4out} (100.0\%), \texttt{unrolled\_dot64}
(95.3\%), and \texttt{unrolled\_dot32} (87.5\%), while
\texttt{sliding\_correlation}, \texttt{multi\_accumulator}, and the
grouped-correlation kernels expose smaller but still measurable exact-free
subsets. These results match the intended use model of the hybrid path:
\textsc{mule} is beneficial only where existing dependence slack or
independent work can fully absorb its additional latency.

The benchmark results should therefore be interpreted conservatively.
Exact-free multiplies are the robust subset whose extra latency remains
fully hidden even after composing the accepted rewrites inside the full
benchmark. A larger middle class of operations exhibits only partial
hiding: these sites satisfy the local trace-based slack heuristic, but
they do not all survive the benchmark-level exact-cycle acceptance test,
typically because individually promising rewrites interfere through
slot-0 issue restrictions, single-outstanding behaviour, or overlapping
dependence chains. The remaining hard-unhidable multiplies are not even
heuristic candidates and therefore correspond to recurrence-dominated
regions such as \texttt{matmul}, \texttt{fir}, \texttt{dotprod}, and
\texttt{horner}. Over the verified corpus, 2,385 of 14,297 dynamic
multiplies (16.7\%) are exact-free, 3,650 (25.5\%) are only partially
hidable, and 8,262 (57.8\%) are hard-unhidable. The main takeaway is
that local slack is common enough to motivate the hybrid path, but only
a smaller composable subset can be harvested without benchmark-level
timing loss.

\paragraph{Mechanics of partial hiding.}
A dynamic multiply site is classified as a \emph{partial-hide candidate}
when the trace-based latency-gap heuristic accepts it: the first
dynamic consumer of the multiply result retires at least
$\Delta L + 1 = 3$ instructions after the multiply itself, providing
enough in-flight work to cover \textsc{mule}'s two extra cycles under
the single-site assumption.  The site fails to become \emph{exact-free},
however, when the full benchmark re-simulation with that PC patched to
\texttt{mule} does not preserve the total cycle count.

Three distinct failure mechanisms account for the gap between the
heuristic acceptance set (6,035 dynamic multiplies, 42.2\% of the
plain-\texttt{mul} pool) and the confirmed exact-free set (2,385,
16.7\%):

\begin{enumerate}
  \item \textbf{Single-outstanding serialisation.}  The \textsc{mule}
  backend accepts at most one operation at a time.  If the loop body
  issues a second \texttt{mule} before the first has completed, the
  issue stage stalls until the backend becomes free.  A site that looks
  slack-sufficient in isolation can become the bottleneck once an
  earlier site in the same iteration also uses \textsc{mule}, because
  the two operations now compete for the single-entry queue.
  A concrete trace example appears in \texttt{multi\_accumulator},
  whose hot loop emits a dense burst of multiplies inside one unrolled
  iteration: \texttt{0x101f2: mul}, \texttt{0x10220: mul},
  \texttt{0x1022e: mul}, \texttt{0x10234: mul}, and subsequent multiply
  sites continue in the same short window, interleaved only with
  accumulator updates.  Under \texttt{mule}, these operations all target
  the same single-entry backend before the earlier one has fully
  drained, so some of the apparent local slack is consumed by queueing
  rather than useful overlap.  This is consistent with
  Table~\ref{tab:benchmark-hiding-breakdown}, where
  \texttt{multi\_accumulator} retains 80 exact-free dynamic multiplies
  but leaves 48 in the partial-hide class.

  \item \textbf{Interference through dependence chains.}  The
  latency-gap filter inspects only the first consumer of each
  multiply.  When an accepted site's result feeds a second-level
  consumer that itself feeds a third, the extra latency can propagate
  and stall a later instruction that was not visible during single-site
  filtering.  \texttt{extended\_estrin} (66.7\% partial) is the most
  affected: its dependence graph is a tree of nested products and
  partial sums, so patching a multiply high in the tree lengthens the
  critical path for all downstream additions, erasing the local slack.
  The trace shows this directly: \texttt{0x10224: mulw a4,a6,a5} feeds
  \texttt{0x1022c: addw a4,a4,a7}, which is then reused by
  \texttt{0x1023c: mulw a4,a4,s9} and folded again at
  \texttt{0x10248: addw s7,s7,a4}.  Other products also re-enter later
  multiply stages, for example \texttt{0x10254: mulw a4,s11,s9} and
  \texttt{0x10266: mulw s10,s10,s9}.  This is why the benchmark can show
  a wide 2-to-25-instruction distance budget in the trace study yet
  still leave 1,536 dynamic multiplies in the partial-hide class once
  those locally safe edges are composed inside the full arithmetic tree.

  \item \textbf{Slot-pressure and dual-issue disruption.}  The
  secondary issue slot (pipe1) can absorb a \texttt{mule} alongside an
  ALU or load/store instruction in the same cycle.  When the scheduler
  cannot pair the \texttt{mule} with an independent instruction---either
  because the instruction window is narrow or because a scoreboard
  hazard forces a bubble---the \texttt{mule} falls back to the primary
  slot and its extra latency is no longer hidden.  This effect is
  visible in \texttt{rdr\_like\_corr}: 50\% of candidates are partial,
  and the all-\texttt{mule} overhead (+27.8\%) is far higher than the
  multiply density (14.1\%) would suggest, indicating frequent
  slot-pressure stalls.  In the extracted hot loop, each body packs four
  multiplies and four accumulator updates into a short issue window:
  \texttt{0x10244: mulw a0,a0,s8}, \texttt{0x1024c: mulw a1,a1,s7},
  \texttt{0x10250: addw t1,a0,t1}, \texttt{0x10254: mulw a2,a2,s6},
  \texttt{0x10258: addw a6,a1,a6}, \texttt{0x1025c: mulw a4,a4,t5}, and
  the final updates at \texttt{0x10260} and \texttt{0x10264}.  The trace
  study reports only 3--4 instructions of immediate mul-to-add spacing
  for these local edges, so there is enough theoretical overlap to mark
  many sites as candidates but not enough consistent pipe1 pairing
  opportunity to hide every converted multiply on the real core.
\end{enumerate}

The practical implication is that a compiler selecting sites for
\texttt{mule} substitution should account for \emph{composition
effects}, not just per-site slack.  Greedy single-site analysis
over-estimates the benefit by 53\% relative to the confirmed
exact-free pool ($6{,}035 - 2{,}385 = 3{,}650$ sites that appear safe
but are not).  A conservative strategy---accepting only sites whose
slack exceeds the worst-case serialisation window, or limiting the
density of \texttt{mule} sites per loop iteration---would bring the
compile-time prediction closer to the benchmark-validated result.

%% ================================================================
\subsection{Break-Even Analysis}\label{sec:eval:breakeven}

For workloads where multiply density is low (i.e., most cycles do
\emph{not} issue a multiply to \textsc{mule}), the hybrid design saves
power simply by having a smaller, lower-leakage second multiplier.
Conversely, for sustained dual-issue multiply workloads, the baseline
$2\!\times\!$MUL is faster. The crossover multiply density $f_m$
(fraction of cycles issuing an efficient-path multiply) at which the hybrid
breaks even on total energy is

\begin{equation}\label{eq:breakeven}
  f_m < \frac{P_{\mathit{mul2}} - P_{\mathit{mule}}}
             {e_{\mathit{cyc}} / T_{\mathit{clk}}}
      = \frac{1.017 - 0.155}{82.70 / 10}
      = \frac{0.862}{8.27}
      \approx 10.4\%.
\end{equation}

\noindent
For workloads where fewer than $\sim$10\% of cycles contain an
efficient-path multiply---which covers the vast majority of
general-purpose and embedded applications---the hybrid configuration
is more energy-efficient. Only sustained, throughput-critical dual-issue
multiply workloads favour the homogeneous baseline.

\begin{table*}[t]
\centering
\caption{Smart-compiler benefit classes derived from the full/partial hiding breakdown. The compiler is assumed to route only benchmark-validated beneficial sites, so exact-free conversions from Table~\ref{tab:benchmark-hiding-breakdown} are selected and partially hidden candidates are rejected. Because every selected conversion preserves total cycle count exactly, any selected benchmark is both energy- and EDP-beneficial under Equations~\ref{eq:net} and \ref{eq:edpratio}.}
\label{tab:benchmark-benefit-classes}
\small
\begin{tabular}{l c r p{0.55\textwidth}}
\hline
Class & Benchmarks & Routed dyn. mul & Benchmarks \\
\hline
EDP benefit & 14 & 2,385 & \texttt{complex\_bank}, \texttt{estrin\_2level}, \texttt{estrin\_poly}, \texttt{extended\_estrin}, \texttt{filter\_bank\_4out}, \texttt{grouped\_correlation\_bank}, \texttt{multi\_accumulator}, \texttt{rdr\_like\_corr}, \texttt{reduction\_tree\_mul}, \texttt{sliding\_correlation}, \texttt{unrolled\_dot8}, \texttt{unrolled\_dot16}, \texttt{unrolled\_dot32}, \texttt{unrolled\_dot64} \\
Energy-only benefit & 0 & 0 & none \\
No benefit & 10 & 0 & \texttt{conv2d}, \texttt{delayed\_reduction}, \texttt{dotprod}, \texttt{fir}, \texttt{fir\_multi}, \texttt{hello}, \texttt{horner}, \texttt{matmul}, \texttt{outer\_product}, \texttt{unrolled\_dot} \\
\hline
\end{tabular}
\end{table*}

Table~\ref{tab:benchmark-benefit-classes} answers the selective-policy
question directly. If the compiler routes only benchmark-validated
beneficial sites, then all 2,385 selected dynamic conversions come from
the exact-free subset and therefore preserve cycle count by
construction. Under that policy, 14 of the 24 kernels become both
energy- and EDP-beneficial, the energy-only class is empty, and the
remaining 10 kernels are left unchanged.

The zero-valued entries in Table~\ref{tab:benchmark-benefit-classes}
are a direct consequence of the hybrid binary threshold. Because
$\Delta e < e_{\mathit{cyc}}$, any routed site that exposes even one
stall cycle is already energy-negative, and therefore cannot populate an
energy-only class. For the same reason, benchmarks listed in the
no-benefit row show zero routed dynamic multiplies: either they contain
no executed plain \texttt{mul} instructions at all, or every locally
promising candidate is rejected once the benchmark-level validation
observes a non-zero cycle penalty.

\begin{table}[t]
\centering
\caption{Per-benchmark estimated energy savings for the 14
EDP-benefiting kernels under the selective policy. Percentages are
relative to each benchmark's all-\texttt{mul} baseline energy
$E_0 = C_{\texttt{mul}} \cdot e_{\mathit{cyc}}$ with
$e_{\mathit{cyc}} = 82.70$\,pJ. The two saving columns differ only in
the per-conversion energy delta: $\Delta e = 52.85$\,pJ from
Table~\ref{tab:params} and $\Delta e_{\mathit{EPI}} = 16.01$\,pJ from
Tables~\ref{tab:params} and \ref{tab:epi_edp}. The multiplication-energy
columns use the in-core total-EPI model: baseline multiplication energy
share is $N_{\texttt{mul}} \cdot 23.76\,\mathrm{pJ} / E_0$, and
validated multiplication-energy saving is
$N_{\texttt{free}} \cdot 16.01\,\mathrm{pJ} / (N_{\texttt{mul}} \cdot 23.76\,\mathrm{pJ})$.}
\label{tab:benchmark-energy-savings}
\small
\resizebox{\columnwidth}{!}{%
\begin{tabular}{l r r r r r}
\hline
Benchmark & Exact-free dyn. mul & Saving (\%, $\Delta e$) & Saving (\%, $\Delta e_{\mathit{EPI}}$) & Mul-energy share & Mul-energy saving \\
\hline
\texttt{estrin\_2level} & 256 & 1.36\% & 0.41\% & 1.53\% & 26.95\% \\
\texttt{rdr\_like\_corr} & 448 & 1.23\% & 0.37\% & 2.20\% & 16.85\% \\
\texttt{extended\_estrin} & 256 & 1.12\% & 0.34\% & 4.55\% & 7.49\% \\
\texttt{estrin\_poly} & 256 & 1.11\% & 0.34\% & 1.25\% & 26.95\% \\
\texttt{sliding\_correlation} & 248 & 0.77\% & 0.23\% & 0.87\% & 26.95\% \\
\texttt{unrolled\_dot64} & 122 & 0.45\% & 0.13\% & 0.21\% & 64.22\% \\
\texttt{unrolled\_dot32} & 112 & 0.42\% & 0.13\% & 0.22\% & 58.96\% \\
\texttt{complex\_bank} & 192 & 0.39\% & 0.12\% & 0.47\% & 25.27\% \\
\texttt{grouped\_correlation\_bank} & 90 & 0.36\% & 0.11\% & 0.43\% & 25.27\% \\
\texttt{unrolled\_dot16} & 88 & 0.36\% & 0.11\% & 0.23\% & 46.33\% \\
\texttt{multi\_accumulator} & 80 & 0.34\% & 0.10\% & 0.24\% & 42.11\% \\
\texttt{unrolled\_dot8} & 80 & 0.34\% & 0.10\% & 0.25\% & 42.11\% \\
\texttt{filter\_bank\_4out} & 125 & 0.16\% & 0.05\% & 0.07\% & 67.38\% \\
\texttt{reduction\_tree\_mul} & 32 & 0.15\% & 0.05\% & 0.27\% & 16.85\% \\
\hline
\end{tabular}
}
\end{table}

Table~\ref{tab:benchmark-energy-savings} makes the benchmark-level
scale concrete. Under the hybrid-model conversion delta
$\Delta e = 52.85$\,pJ, the 14 EDP-benefiting kernels span an estimated
energy reduction from 0.15\% (\texttt{reduction\_tree\_mul}) to 1.36\%
(\texttt{estrin\_2level}) of their own all-\texttt{mul} baseline
energy. Using the more conservative in-core EPI gap
$\Delta e_{\mathit{EPI}} = 16.01$\,pJ preserves the same ordering but
compresses the range to 0.05\%--0.41\%. Under that same EPI model, the
baseline multiplication-energy share of these kernels ranges from
0.07\% (\texttt{filter\_bank\_4out}) to 4.55\%
(\texttt{extended\_estrin}), while the validated selective policy saves
between 7.49\% and 67.38\% of the multiplication energy within each
benchmark, depending on how much of the dynamic plain-\texttt{mul} pool
falls into the exact-free subset.

The partial-hide class is still informative, but now only as a
rejection set. Across the full corpus, 3,650 dynamic multiplies appear
locally promising yet fail benchmark-level validation once composition
effects are considered. Two kernels, \texttt{delayed\_reduction} and
\texttt{unrolled\_dot}, fall into a particularly clear partial-only
category: they contain local slack candidates but no benchmark-validated
beneficial \texttt{mule} routes. Aggregated over all 24 kernels, the
smart-compiler policy therefore saves 126.0\,nJ relative to the
all-\texttt{mul} baseline, or about 0.32\% of the baseline benchmark
energy, under the hybrid-model delta. Under the local in-core EPI gap,
the same 2,385 validated conversions correspond to 38.2\,nJ, or about
0.10\% of the same baseline energy, while still introducing zero
benchmark slowdown.

%% ================================================================
\subsection{The Role of Loop Unrolling}\label{sec:eval:unrolling}

The verified dot-product family highlights the compiler side of the
co-design problem. The baseline \texttt{dotprod} kernel exposes no
exact-free rewrites, and the lightly unrolled \texttt{unrolled\_dot}
variant still slows down badly under the all-\texttt{mule} experiment.
By contrast, the wider unrolled variants recover substantial hidden
slack: \texttt{unrolled\_dot8} and \texttt{unrolled\_dot16} admit
62.5\% and 68.8\% exact-free rewrites, respectively, while
\texttt{unrolled\_dot32} and
\texttt{unrolled\_dot64} reach 87.5\% and 95.3\%, respectively, with
all selective-rewrite runs preserving cycle count exactly.

This progression underscores a key insight: \emph{the energy benefit of
slack-aware functional-unit selection is not solely a hardware
property---it depends on compiler transformations that reshape the
dependence graph}. Loop unrolling, accumulator replication, software
pipelining, and instruction scheduling all enlarge the window in which
a slower but cheaper execution path can absorb a multiply without
changing end-to-end timing.

%% ================================================================
\subsection{Discussion}\label{sec:eval:discussion}

The combined evidence from Section~\ref{sec:eval:switching2026} and
Section~\ref{sec:eval:bench} supports a conditional claim: the Hybrid
architecture can outperform switched 2$\times$MUL on selected workloads,
but switching itself is beneficial only when added iterative latency is
sufficiently hidden after benchmark-level composition effects.

\paragraph{Leakage at advanced nodes.}
At 45\,nm, leakage accounts for 6.9--14.8\% of the standalone
multiplier power. At advanced nodes (16\,nm, 7\,nm), leakage rises to
30--50\% of total power. Because \textsc{mule}'s area is 52.7\%
smaller, its leakage advantage compounds: at a hypothetical 50\%
leakage fraction, the multiplier area reduction alone would save
$\sim$2.0\% of core power, on top of the dynamic savings.

\paragraph{VCD annotation coverage.}
Our power numbers are based on $\sim$10--11\% VCD annotation of flop
outputs, with the remainder estimated at a default switching activity
of 0.2. This is conservative: actual switching is often lower than 0.2
for idle logic, so the reported power difference likely
\emph{understates} \textsc{mule}'s advantage. Higher-coverage
annotation (e.g., gate-level simulation) would tighten the bounds.

\paragraph{Threats to validity.}
The evaluation uses a single technology node (45\,nm) and a single
operating point (100\,MHz, 1.1\,V). While the qualitative conclusions
(binary threshold, area--energy tradeoff) are technology-independent,
the quantitative margins will shift at different nodes, voltages, and
frequencies. Additionally, the 24 benchmark kernels are still synthetic
microbenchmarks rather than a full application suite; broader workloads
(e.g., SPEC, Embench) are needed to characterise the population
distribution of multiply slack across real applications.